% -*- coding: utf-8 -*-
\documentclass[aspectratio=169]{beamer}

\usetheme{UFS}

\usepackage[utf8]{inputenc}
\usepackage[brazilian]{babel}
\usepackage{amsmath,amssymb}
\usepackage{graphicx}
\usepackage{booktabs}
\usepackage{listings}
\usepackage{xcolor}
\usepackage{tikz}
\usetikzlibrary{positioning, shapes.multipart, arrows.meta, shadows, fit, chains, decorations.pathreplacing, shapes.geometric, calc}

\definecolor{codegreen}{rgb}{0,0.6,0}
\definecolor{codegray}{rgb}{0.5,0.5,0.5}
\definecolor{codepurple}{rgb}{0.58,0,0.82}
\definecolor{backcolour}{rgb}{0.95,0.95,0.92}

\lstdefinestyle{mystyle}{
    backgroundcolor=\color{backcolour},
    commentstyle=\color{codegreen},
    keywordstyle=\color{magenta},
    numberstyle=\tiny\color{codegray},
    stringstyle=\color{codepurple},
    breaklines=true,
    captionpos=b,
    keepspaces=true,
    numbers=left,
    numbersep=5pt,
    showstringspaces=false,
    tabsize=2,
    frame=single,
    rulecolor=\color{gray},
    basicstyle=\ttfamily\scriptsize,
}
\lstset{style=mystyle,language=C}

\title{Hashing: rehashing e complexidade amortizada}
\subtitle{Aplicações: conjuntos, dicionários, contagem de frequências}
\author[André Y. Kashiwabara]{Prof. Dr. André Yoshiaki Kashiwabara}
\institute{UFS -- COMP0497 Estruturas de Dados}
\date{\today}

\begin{document}

\begin{frame}[plain,noframenumbering]
    \titlepage
\end{frame}

\begin{frame}{Roteiro da aula (150 min)}
\tableofcontents
\end{frame}

\section{Recapitulação}

\begin{frame}{O que já vimos sobre tabelas hash}
\begin{itemize}
    \item Função hash $h: U \to \{0,\dots,m-1\}$.
    \item Tratamento de colisões: \textbf{encadeamento} e \textbf{endereçamento aberto}
    (sondagem linear, quadrática, hash duplo).
    \item Fator de carga $\alpha = n/m$.
    \item Custo \emph{esperado} de busca:
    \begin{itemize}
        \item Encadeamento: $\Theta(1 + \alpha)$.
        \item Endereçamento aberto: $\Theta(1/(1-\alpha))$ (uniform hashing).
    \end{itemize}
    \item Quando $\alpha$ cresce, o desempenho \textbf{degrada}.
\end{itemize}
\end{frame}

\begin{frame}{Pergunta motivadora}
\begin{block}{Cenário}
Inserimos um milhão de chaves em uma tabela de tamanho 1024. Qual o problema?
\end{block}
\bigskip
\begin{itemize}
    \item $\alpha \approx 977$. Buscas lineares em listas enormes.
    \item Encadeamento: cada bucket vira uma lista de $\sim 1000$ elementos.
    \item Endereçamento aberto: tabela cheia $\Rightarrow$ inserção \emph{impossível}.
\end{itemize}
\bigskip
\textbf{Solução:} crescer a tabela dinamicamente. Mas como, sem perder a localização
das chaves já inseridas?
\end{frame}

\section{Rehashing}

\begin{frame}{O que é \emph{rehashing}}
\begin{itemize}
    \item Quando $\alpha$ ultrapassa um \textbf{limiar} $\alpha_{\max}$, criamos
    uma nova tabela $T'$ com tamanho $m' > m$.
    \item Para cada chave $k$ em $T$:
    \begin{itemize}
        \item Recalcular $h'(k) \bmod m'$.
        \item Inserir em $T'$.
    \end{itemize}
    \item Liberar a tabela antiga.
\end{itemize}
\bigskip
\textbf{Pontos críticos:}
\begin{itemize}
    \item Não basta copiar: as posições mudam porque $m$ mudou.
    \item Custo de uma operação que dispara rehash: $\Theta(n)$.
    \item Operações ``normais'': $\Theta(1)$.
    \item Como conciliar?
\end{itemize}
\end{frame}

\begin{frame}{Estratégias de crescimento}
\begin{itemize}
    \item \textbf{Crescimento aditivo} ($m' = m + c$): péssimo. Discutiremos por quê.
    \item \textbf{Crescimento multiplicativo} ($m' = 2m$, ou fator $1{,}5\times$): padrão
    de mercado.
    \item \textbf{Tamanho primo:} útil para reduzir colisões patológicas; escolhe-se
    o próximo primo após $2m$.
    \item \textbf{Encolhimento:} ao remover muitas chaves, podemos reduzir $m$ se
    $\alpha$ cair abaixo de outro limiar (com cuidado: ver \emph{thrashing}).
\end{itemize}
\end{frame}

\begin{frame}{Limiar e \emph{thrashing}}
\begin{itemize}
    \item Limiar de crescimento típico: $\alpha_{\text{up}} = 0{,}75$ (Java HashMap),
    $0{,}66$ (Python dict), $0{,}5$ (endereçamento aberto agressivo).
    \item Se reduzirmos quando $\alpha < \alpha_{\text{up}}/2$, evitamos \emph{thrashing}:
    sequência alternada de \emph{insert/remove} disparando rehash a cada operação.
    \item Use sempre uma \textbf{histerese} entre os limiares.
\end{itemize}
\end{frame}

\section{Análise amortizada}

\begin{frame}{Por que ``amortizada''?}
\begin{itemize}
    \item Operações isoladas variam:
    \begin{itemize}
        \item Maioria: $\Theta(1)$.
        \item Algumas: $\Theta(n)$ (rehash).
    \end{itemize}
    \item No \emph{pior caso} por operação, $\Theta(n)$.
    \item No \emph{caso médio sobre uma sequência} de operações, queremos algo melhor.
    \item \textbf{Análise amortizada:} custo \emph{médio por operação} sobre uma
    sequência de $n$ operações no pior caso.
\end{itemize}
\end{frame}

\begin{frame}{Três técnicas (Cormen et al., cap. 17)}
\begin{enumerate}
    \item \textbf{Análise agregada:} somar custos totais, dividir por $n$.
    \item \textbf{Método contábil (\emph{accounting}):} cobrar mais por algumas
    operações para ``pagar adiantado'' o custo de outras.
    \item \textbf{Método potencial:} associar a cada estado uma função $\Phi$,
    e escrever o custo amortizado como $\hat c_i = c_i + \Phi_i - \Phi_{i-1}$.
\end{enumerate}
\bigskip
Vamos aplicar a primeira ao crescimento de tabelas.
\end{frame}

\begin{frame}{Crescimento dobrando: análise agregada}
\begin{itemize}
    \item Começamos com $m_0 = 1$. Inserimos $n$ chaves seguidas.
    \item A tabela dobra quando $\alpha$ atinge 1: tamanhos $1, 2, 4, 8, \dots$.
    \item Quantos rehashes? $\lceil \log_2 n \rceil$.
    \item Custo de cada rehash: copiar todas as chaves atuais.
    \item Soma:
    \[
    \sum_{i=0}^{\lfloor \log_2 n \rfloor} 2^i = 2^{\lfloor \log_2 n \rfloor + 1} - 1 < 2n.
    \]
    \item Custo total das inserções: $n + 2n = 3n$.
    \item \textbf{Custo amortizado:} $3n/n = O(1)$ por inserção.
\end{itemize}
\end{frame}

\begin{frame}{Por que crescer aditivamente é desastroso}
\begin{itemize}
    \item Suponha $m' = m + c$, $c$ constante.
    \item Inserir $n$ chaves $\Rightarrow$ $n/c$ rehashes.
    \item Custos cumulativos: $c, 2c, 3c, \dots, n$.
    \item Soma: $\Theta(n^2/c) = \Theta(n^2)$.
    \item \textbf{Custo amortizado:} $\Theta(n)$ por inserção.
\end{itemize}
\bigskip
\textbf{Moral:} dobrar (ou multiplicar por constante $> 1$) é o que torna
\textbf{rehashing} viável.
\end{frame}

\begin{frame}{Método potencial aplicado a tabelas dinâmicas}
\begin{itemize}
    \item Seja $T_i$ o estado da tabela após operação $i$, com $n_i$ elementos e
    tamanho $m_i$.
    \item Função potencial: $\Phi(T) = 2 n - m$ (quando $n \geq m/2$).
    \item Custo amortizado de inserção:
    \begin{itemize}
        \item Sem rehash: $\hat c = 1 + \Delta \Phi = 1 + 2 = 3$.
        \item Com rehash ($n_{i-1} = m_{i-1}$, $m_i = 2 m_{i-1}$):
        \[
        \hat c = (n_{i-1} + 1) + \Delta \Phi = (n+1) + (2 - n) = 3.
        \]
    \end{itemize}
    \item \textbf{Conclusão:} cada inserção custa, amortizadamente, $O(1)$.
\end{itemize}
\end{frame}

\begin{frame}{Visualizando o custo}
\begin{center}
\begin{tikzpicture}[scale=0.85]
\draw[->] (0,0) -- (12,0) node[right] {$n$};
\draw[->] (0,0) -- (0,4) node[above] {custo};
\foreach \i/\h in {1/0.3,2/0.3,3/0.6,4/0.3,5/0.9,6/0.3,7/0.3,8/0.3,9/1.5,10/0.3,11/0.3} {
    \draw[fill=blue!30] (\i-0.4,0) rectangle (\i+0.4,\h);
}
\draw[red, thick, dashed] (0.5, 0.5) -- (11.5, 0.5) node[right] {amortizado};
\end{tikzpicture}
\end{center}
\begin{itemize}
    \item Picos em $n = 2, 4, 8, \dots$ — momentos de rehash.
    \item Linha tracejada: custo amortizado constante.
    \item ``Pague três por inserção'': dois ficam de \emph{crédito} para futuro rehash.
\end{itemize}
\end{frame}

\begin{frame}{Inserção + remoção: ainda $O(1)$ amortizado?}
\begin{itemize}
    \item Se reduzirmos quando $\alpha < 1/2$: pode ocorrer \emph{thrashing}.
    \item Solução: reduzir só quando $\alpha < 1/4$ (e crescer quando $\alpha = 1$).
    \item Função potencial híbrida:
    \[
    \Phi = \begin{cases}
        2n - m & \text{se } \alpha \geq 1/2 \\
        m/2 - n & \text{se } \alpha < 1/2
    \end{cases}
    \]
    \item Análise (Cormen 17.4) mostra: \textbf{inserção e remoção custam $O(1)$
    amortizado.}
\end{itemize}
\end{frame}

\section{Pior caso vs.\ amortizado}

\begin{frame}{Quando o amortizado não basta}
\begin{itemize}
    \item Sistemas \textbf{tempo real}: latência de uma única operação importa.
    \item Servidores web sob latência cauda (\emph{tail latency}): um rehash de
    $10^7$ chaves pausa o serviço por dezenas de ms.
    \item \textbf{Rehashing incremental:} amortizar o trabalho ao longo de várias
    operações.
\end{itemize}
\end{frame}

\begin{frame}{Rehashing incremental}
\begin{itemize}
    \item Manter \emph{duas} tabelas em paralelo: $T_{\text{antiga}}$ e $T_{\text{nova}}$.
    \item Cada operação:
    \begin{itemize}
        \item Move $k$ chaves de $T_{\text{antiga}}$ para $T_{\text{nova}}$.
        \item Realiza a operação, consultando ambas se necessário.
    \end{itemize}
    \item Inserções vão direto em $T_{\text{nova}}$.
    \item Quando $T_{\text{antiga}}$ fica vazia, é descartada.
    \item \textbf{Custo por operação:} $O(k)$ no pior caso.
    \item Usado em Redis, alguns kernels e bancos de dados em memória.
\end{itemize}
\end{frame}

\section{Aplicações}

\begin{frame}{Conjuntos (\emph{sets})}
\begin{itemize}
    \item Operação esperada $O(1)$: \texttt{add}, \texttt{contains}, \texttt{remove}.
    \item Implementação: tabela hash com chaves apenas (sem valor associado).
    \item Operações de conjunto: união, interseção, diferença.
    \item Aplicações: deduplicação, controle de visitados em grafos, filtros de
    spam (\emph{bloom filters}: hashing probabilístico).
\end{itemize}
\end{frame}

\begin{frame}{Dicionários (\emph{maps})}
\begin{itemize}
    \item Estrutura mais usada em programação moderna.
    \item Java: \texttt{HashMap} (encadeamento $\to$ árvore quando bucket cresce muito,
    desde Java 8).
    \item Python: \texttt{dict} (endereçamento aberto, sondagem perturbada,
    preserva ordem de inserção desde 3.7).
    \item C++: \texttt{std::unordered\_map} (encadeamento).
    \item Aplicações: índices, caches, \emph{memoization}, contagem.
\end{itemize}
\end{frame}

\begin{frame}[fragile]{Contagem de frequências}
\begin{lstlisting}[language=Python]
from collections import defaultdict

def contar(palavras):
    freq = defaultdict(int)
    for p in palavras:
        freq[p] += 1
    return freq
\end{lstlisting}
\begin{itemize}
    \item Tempo total: $\Theta(N)$ esperado, onde $N$ = número de palavras.
    \item Tempo com BST: $\Theta(N \log U)$, onde $U$ = vocabulário.
    \item Hash vence quando ordenação não é necessária.
    \item Aplicação clássica: \emph{word count} em logs, MapReduce, NLP.
\end{itemize}
\end{frame}

\begin{frame}{Outras aplicações}
\begin{itemize}
    \item \textbf{Caches} (LRU): hash + lista duplamente ligada $\Rightarrow$ $O(1)$
    para \texttt{get} e \texttt{put}.
    \item \textbf{Memoization:} cachear resultados de função pura.
    \item \textbf{\emph{Symbol tables}} em compiladores.
    \item \textbf{Detecção de duplicatas} em streams.
    \item \textbf{Hashing criptográfico:} SHA-256, integridade de arquivos
    (uso diferente: resistência a colisão deliberada).
    \item \textbf{Hashing consistente:} balanceamento em sistemas distribuídos
    (caches, sharding).
\end{itemize}
\end{frame}

\section{Considerações práticas}

\begin{frame}{Escolha da função hash}
\begin{itemize}
    \item Boa função: distribuição quase-uniforme sobre $\{0,\dots,m-1\}$.
    \item Strings: \emph{polynomial rolling}, FNV, MurmurHash, SipHash.
    \item Inteiros: multiplicativo de Knuth ($h(k) = \lfloor m\,(k\phi \bmod 1) \rfloor$).
    \item Em ambientes adversariais (entrada controlada por atacante), use hash
    com \textbf{semente aleatória} (ex.\ SipHash em Python, Rust).
    \item \textbf{Evite:} divisão por $m$ não-primo com chaves correlacionadas.
\end{itemize}
\end{frame}

\begin{frame}{Ataques de colisão}
\begin{itemize}
    \item 2003: ataque a Perl, PHP, Java; sequências escolhidas geram colisões em
    cascata.
    \item Sintoma: requisição HTTP com 1000 parâmetros causa \emph{denial of service}.
    \item Mitigação: \textbf{hash com chave secreta} (SipHash, randomização de seed
    por execução).
    \item Java 8 mitiga via \emph{treeify}: bucket vira árvore vermelho-preta quando
    excede 8 elementos.
\end{itemize}
\end{frame}

\begin{frame}{Resumo}
\begin{itemize}
    \item \textbf{Rehashing} permite tabelas hash dinâmicas sem perder $O(1)$
    esperado por operação.
    \item \textbf{Crescimento multiplicativo} é essencial: aditivo dá $\Theta(n)$
    amortizado.
    \item \textbf{Análise amortizada} formaliza o custo médio sobre sequências:
    agregada, contábil, potencial.
    \item Para sistemas \emph{soft real-time}, prefira \textbf{rehashing incremental}.
    \item Aplicações dominam programação moderna: conjuntos, dicionários, contagem,
    caches.
\end{itemize}
\end{frame}

\begin{frame}{Para a próxima aula}
\begin{itemize}
    \item Implementar uma tabela hash com rehashing automático em C.
    \item Medir tempo por operação em uma sequência de $10^6$ inserções.
    \item Comparar com \texttt{std::unordered\_map}.
    \item Refletir: como você projetaria um \emph{dict} para um sistema de
    trading com latência cauda $< 1$ ms?
\end{itemize}
\centering\Large Obrigado!
\end{frame}

\end{document}
